В ходе выполнения курсовой работы была рассмотрена актуальная научно-техническая задача, связанная с повышением точности измерения углового положения осей телескопических монтировок.

Значимость данной задачи определяется возрастающей потребностью в автоматизации астрономических наблюдений, проведении массовых обзоров неба и создании высокоточных систем наведения, способных функционировать в составе современных роботизированных измерительных комплексов. Особую практическую важность работа приобретает в условиях дефицита отечественных прецизионных энкодеров, пригодных для применения в астрономическом приборостроении.

В процессе исследования был выполнен анализ существующих типов абсолютных датчиков угла, что позволило обосновать выбранные технические решения и сформировать требования к разрабатываемому прототипу. Был создан и экспериментально исследован макет абсолютного энкодера, проведены его отладка и серия испытаний. Полученные результаты позволили выявить как случайные, так и систематические составляющие погрешности измерений, а также оценить влияние конструктивных и алгоритмических факторов на точность работы устройства.

Для обработки сигнала были применены методы математического и программного анализа, подтвердившие возможность повышения качества измерений за счет программно-алгоритмической коррекции. Таким образом, поставленные в курсовой работе цели были достигнуты, а разработанный прототип продемонстрировал перспективность дальнейшего совершенствования. Результаты работы могут быть использованы при создании отечественных высокоточных систем углового позиционирования для астрономических телескопов и иных измерительных установок.